\documentclass[10pt]{article}
\usepackage[T1]{fontenc}
\usepackage[utf8]{inputenc}
\usepackage[french]{babel}
\usepackage{amsmath,amssymb,mathrsfs,mathpazo}
\usepackage{geometry}
\geometry{a4paper,margin=2cm}
\usepackage{graphicx,array,tabularx,tabularray,colortbl,multirow,multicol,booktabs,enumitem,pifont,ulem,fancybox,diagbox,cancel}
\usepackage{pstricks,pstricks-add,pst-plot,pst-tree,pst-node,pst-text,pst-func,pst-math,pst-bspline,pst-3dplot,pst-eucl,pst-solides3d}
\let\pscalculate\undefined
\usepackage{pst-bezier}
\usepackage{tikz,pgfplots,tkz-tab}
\usetikzlibrary{arrows,arrows.meta,patterns,calc,positioning,trees,decorations.markings,shapes.geometric,shapes.misc,shapes.arrows,intersections,decorations.pathmorphing,decorations.pathreplacing,backgrounds,shadows,babel}
\pgfplotsset{compat=1.18}
\usepgfplotslibrary{fillbetween}
\usepackage[np]{numprint}
\usepackage{listings,makecell,pas-tableur,twoopt}
\let\barre\undefined
\usepackage[pstricks]{tablvar}
\definecolor{darkgreen}{RGB}{0,100,0}
\providecommand{\red}{\color{red}}
\providecommand{\blue}{\color{blue}}
\providecommand{\green}{\color{green}}
\providecommand{\black}{\color{black}}
\providecommand{\white}{\color{white}}
\providecommand{\gray}{\color{gray}}
\providecommand{\cyan}{\color{cyan}}
\providecommand{\magenta}{\color{magenta}}
\providecommand{\yellow}{\color{yellow}}
\providecommand{\violet}{\color{violet}}
\providecommand{\orange}{\color{orange}}
\providecommand{\euro}{\texteuro}
\providecommand{\EUR}{\texteuro}
\providecommand{\ieme}{\textsuperscript{e}}
\providecommand{\degres}{\ensuremath{^\circ}}
\providecommand{\degree}{\ensuremath{^\circ}}
\providecommand{\vv}[1]{\overrightarrow{#1}}
\providecommand{\Coord}[1]{#1}
\providecommand{\decosix}{}
\providecommand{\starredbullet}{\ensuremath{\bullet}}
\pagestyle{empty}
\setlength{\parindent}{0pt}
\usepackage{adjustbox}\DeclareUnicodeCharacter{2236}{:}\DeclareUnicodeCharacter{2190}{\ensuremath{\leftarrow}}
\providecommand{\R}{\mathbb{R}}
\providecommand{\N}{\mathbb{N}}
\providecommand{\Z}{\mathbb{Z}}
\providecommand{\Q}{\mathbb{Q}}
\providecommand{\D}{\mathbb{D}}
\providecommand{\C}{\mathbb{C}}
\providecommand{\e}{{\mathrm{e}}}
\providecommand{\vect}[1]{\overrightarrow{#1}}
\providecommand{\vectt}[1]{\overrightarrow{\mathrm{#1}}}
\providecommand{\vvt}[1]{\overrightarrow{\mathrm{#1}}}
\providecommand{\np}[2][]{{#2~#1}}
\providecommand{\vv}[1]{{\overrightarrow{#1}}}
\renewcommand{\d}{{\mathrm{d}}}
\providecommand{\degres}{{^\circ}}
\providecommand{\degre}{{^\circ}}
\providecommand{\blue}{\color{blue}}
\providecommand{\red}{\color{red}}
\providecommand{\enskip}{\quad}
\providecommand{\llbracket}{\lbrack\!\lbrack}
\providecommand{\rrbracket}{\rbrack\!\rbrack}
\providecommand{\euro}{{\text{€}}}
\providecommand{\no}{{\text{n°}}}
\providecommand{\interval}[2]{{[#1;#2]}}
\providecommand{\intervFO}[2]{{[#1;#2[}}
\providecommand{\dotfill}{\ldots}

\providecommand{\ds}{\displaystyle}
\providecommand{\textding}[1]{#1}
\providecommand{\Oij}{\(\mathrm{O};\vec{\imath},\vec{\jmath}\)}
\providecommand{\Oijk}{\(\mathrm{O};\vec{\imath},\vec{\jmath},\vec{k}\)}
\providecommand{\Ouv}{\(\mathrm{O};\vec{u},\vec{v}\)}
\providecommand{\intervFF}[2]{\ensuremath{[#1;#2]}}
\renewcommand{\euro}{\ifmmode\text{\texteuro}\else\texteuro\fi}
\DeclareUnicodeCharacter{00B0}{\ifmmode^{\circ}\else\textdegree\fi}
\DeclareUnicodeCharacter{20AC}{\ifmmode\text{\texteuro}\else\texteuro\fi}
% Énoncés : Éducation nationale. Transcription et corrigés : APMEP, Denis Vergès et contributeurs. Découpage et classement : CoopMaths, Sébastien Lozano, Éric Elter et contributeurs. Adaptation de la présentation : Maths Explorer.
% CC BY-SA : https://coopmaths.fr/www/about/
% https://coopmaths.fr/alea/static/bac/2023/tex/bac_2023_08_sujet1_caledonie_4.tex
\begin{document}
\noindent\textit{APMEP / CoopMaths. CC BY-SA. Présentation adaptée par Maths Explorer.}\par\medskip
\medskip

On considère la fonction $f$ définie pour tout réel $x$ de l'intervalle $]0~;~+\infty[$  par :
\[f(x) = 5x^2 + 2x - 2x^2\ln (x).\]

On note $\mathcal{C}_f$ la courbe représentative de $f$ dans un repère orthogonal du plan.

On admet que $f$ est deux fois dérivable sur l'intervalle $]0~;~+\infty[$.

On note $f'$ sa dérivée et $f''$ sa dérivée seconde.

\medskip

\begin{enumerate}
\item 
	\begin{enumerate}
		\item Démontrer que la limite de la fonction $f$ en $0$ est égale à $0$.
		\item Déterminer la limite de la fonction $f$ en $+ \infty$.
	\end{enumerate}
\item Déterminer $f'(x)$ pour tout réel $x$ de l'intervalle $]0~;~+\infty[$.
\item 
	\begin{enumerate}
		\item Démontrer que pour tout réel $x$ de l'intervalle $]0~;~+\infty[$
\[f''(x) = 4(1 - \ln (x)).\]
		\item En déduire le plus grand intervalle sur lequel la courbe $\mathcal{C}_f$ est au-dessus de ses tangentes.
		\item Dresser le tableau des variations de la fonction $f'$ sur l'intervalle $]0~;~+\infty[$.

(On admettra que $\begin{gathered}\displaystyle\lim_{\substack{x \to 0\\x>0}} f'(x) = 2\end{gathered}$ et que $\displaystyle\lim_{x \to + \infty} f'(x) = - \infty$.)
	\end{enumerate}	
\item 
	\begin{enumerate}
		\item Montrer que l'équation $f'(x) = 0$ admet dans l'intervalle $]0~;~+\infty[$ une unique solution $\alpha$ dont on donnera un encadrement d'amplitude $10^{-2}$.
		\item En déduire le signe de $f'(x)$ sur l'intervalle $]0~;~+\infty[$ ainsi que le tableau des variations de la fonction $f$ sur l'intervalle $]0~;~+\infty[$.
	\end{enumerate}
\item 
	\begin{enumerate}
		\item En utilisant l'égalité $f'(\alpha) = 0$, démontrer que : 
		
\[\ln (\alpha) = \dfrac{4\alpha + 1}{2\alpha}.\]
 En déduire que $f(\alpha) = \alpha^2 + \alpha$.
		\item En déduire un encadrement d'amplitude $10^{-1}$ du maximum de la fonction $f$.
	\end{enumerate}
\end{enumerate}

\end{document}
